\subsection{ASAP7 Synthesis, Timing, Area, and Controller Power}
\label{sec:asap7-controller-characterization}

We synthesized the clock-gated controller with OpenROAD Flow Scripts using
the predictive ASAP7 7.5-track RVT standard-cell platform.  The synthesis
constraint is a 1.000-ns core clock, matching the 1-GHz ONNXim core frequency,
with 200-ps input and output delays.  The implementation contains the complete
controller RTL described in Sec.~\ref{sec:rtl-clock-gating}, including both
resident descriptors, causal-bound arithmetic, the three-stage certificate
pipeline, the countdown controller, two explicit integrated clock gates, and
two local clock buffers.

Table~\ref{tab:controller-physical-summary} reports the retained
post-synthesis results.  The timing report uses ASAP7 RVT fast-corner NLDM
libraries and an ideal clock network.  At the 1.000-ns target, the worst
reported setup slack is positive by 224.88~ps.  The minimum reported clock
period is 760.46~ps, corresponding to 1.31499~GHz.  We use 1~GHz in all
architectural and energy calculations; the reported maximum frequency is a
synthesis-stage timing result rather than the operating frequency assumed by
the evaluation.

\begin{table}[t]
  \centering
  \caption{Post-synthesis controller implementation in ASAP7.}
  \label{tab:controller-physical-summary}
  \begin{tabular}{lr}
    \toprule
    Metric & Measured value \\
    \midrule
    Target clock period & 1.000 ns \\
    Worst setup slack & $+224.88$ ps \\
    Minimum reported period & 760.46 ps \\
    Reported maximum frequency & 1.31499 GHz \\
    Standard-cell area & 435 $\mu\mathrm{m}^{2}$ \\
    Mapped integrated clock gates & 2 \\
    Mapped local clock buffers & 2 \\
    \bottomrule
  \end{tabular}
\end{table}

The controller area is compared with the structurally extrapolated
128$\times$128 SA area in Sec.~\ref{sec:sa-scaling}.  Using
3.800769~mm$^{2}$ for that extrapolated array gives a controller-to-SA area
ratio of 0.01145\%.  This percentage inherits the array-area extrapolation and
must not be interpreted as a measured placed 128$\times$128 area overhead.

Controller power is evaluated at the ASAP7 RVT TT corner using workload
activity rather than a default toggle rate.  Each architectural controller
trace is replayed through the mapped gate-level netlist, producing a VCD over
the full workload interval after reset.  OpenSTA reports 13,207 annotated pin
activities and zero unannotated functional pins for each accepted sparse-event
power report.  Reset activity is excluded.  Table~\ref{tab:controller-power-example}
shows the exact four-core power values for the end-to-end OPT-125M execution
as a representative example of the procedure.  Energy is not obtained by
averaging these powers: for each core, its measured total power is multiplied
by its own trace duration, and the resulting energies are summed.

\begin{table}[t]
  \centering
  \caption{Workload-annotated controller power for end-to-end OPT-125M.}
  \label{tab:controller-power-example}
  \begin{tabular}{lrrrr}
    \toprule
    Core & Internal & Switching & Leakage & Total \\
         & \multicolumn{4}{c}{$(\mu\mathrm{W})$} \\
    \midrule
    0 & 37.8856 & 4.8548 & 0.4335 & 43.1739 \\
    1 & 38.1459 & 5.0854 & 0.4297 & 43.6611 \\
    2 & 37.8744 & 4.8502 & 0.4337 & 43.1583 \\
    3 & 38.1302 & 5.0742 & 0.4299 & 43.6342 \\
    \bottomrule
  \end{tabular}
\end{table}

These are synthesis-stage, zero-delay activity results.  They include mapped
cell internal, switching, and leakage power, including the explicit clock
gates and local buffers, but exclude routed-interconnect capacitance,
post-route clock-tree effects, and delay-dependent glitches.  Consequently,
they are not signoff-quality post-layout power estimates.

\subsection{HSPICE SA and Header-Switch Characterization}
\label{sec:hspice-characterization}

Standard-cell timing and power analysis do not model the virtual-supply rail
collapse, wake-up inrush, or transition energy of a power-gated domain.  We
therefore characterize the connected SA load with Synopsys HSPICE using the
ASAP7 TT transistor model at 0.7~V and 25$^{\circ}$C.  The connected 1$\times$1
test vehicle is the mapped two-stage weight-stationary PE netlist containing
2,225 standard-cell instances.  ASAP7 xAct3D-extracted RVT cell subcircuits
provide the logic load.  The power switch is modeled as a lumped parallel bank
of RVT PMOS headers; each unit device has $W=81$~nm, $L=20$~nm, and three fins.

The characterized boundary is deliberately narrow.  Before sleep, the PE is
drained to a reset-zero state, its clock is stopped low, and its top-level
inputs are clamped low.  The PE state is non-retentive.  The model excludes
output-isolation cells, a physical sleep-driver implementation, distributed
header placement, and placed power-grid parasitics.  The sleep command uses a
50-ps slew in the connected-block characterization, and control-source energy
is included explicitly in the event energy.

Unit-header DC measurements first relate average active current to candidate
virtual-rail-drop targets.  For a candidate drop $d$, the number of parallel
devices is
\begin{equation}
  M(d)=\left\lceil\frac{P_{\mathrm{active}}/V_{DD}}
  {I_{\mathrm{unit}}(d)}\right\rceil .
  \label{eq:header-sizing}
\end{equation}
The measured unit currents and resulting 1$\times$1 candidates are listed in
Table~\ref{tab:header-sizing-sweep}.  This calculation uses average active
current; it is not a worst-case instantaneous-droop or electromigration
signoff analysis.

\begin{table}[t]
  \centering
  \caption{Current-derived PMOS-header sizing points.}
  \label{tab:header-sizing-sweep}
  \begin{tabular}{rrr}
    \toprule
    Drop target (\%) & Unit ON current ($\mu$A) & Multiplicity \\
    \midrule
    10.0 & 24.8138 & 44 \\
     7.5 & 19.0019 & 57 \\
     5.0 & 12.9062 & 83 \\
     2.5 &  6.5601 & 164 \\
    \bottomrule
  \end{tabular}
\end{table}

At the 5\% current-derived point ($M=83$), the directly connected 1$\times$1
zero-state experiment measures 7.2927~$\mu$s from the sleep-control 50\%
crossing until the virtual rail falls through 10\% of $V_{DD}$, and
477.365~ps from the wake-control 50\% crossing until the rail rises through
90\% of $V_{DD}$.  The measured steady ON and OFF powers are 352.652~nW and
3.8979~nW, respectively.  These measurements characterize that sizing point;
they are not reused as the transition parameters of a differently sized
header.

The final workload energy calculation uses a separate directly simulated
duration sweep with an $M=300$ equivalent header per 1$\times$1 PE.  The
retained artifact establishes that this is the evaluated header size, but it
does not contain a derivation showing that $M=300$ is an unconstrained optimum
or a numerical worst-case active-droop guarantee.  We therefore treat
$M=300$ as an explicit sizing assumption and do not claim it is optimal.

For the $M=300$ sweep, the directly measured ON reference is
351.8253~nW/PE and the steady OFF power is 14.0768~nW/PE, a 95.9989\%
steady-state reduction.  Fifteen sleep-plateau durations from 5~ns to
15~$\mu$s are simulated.  For each duration, the event cost is the sum of
energy drawn from $V_{DD}$ and the absolute energy at the ideal sleep-control
source.  Selected measured points are shown in
Table~\ref{tab:m300-duration-sweep}.

\begin{table}[t]
  \centering
  \caption{Selected points from the direct $M=300$ duration sweep.}
  \label{tab:m300-duration-sweep}
  \begin{tabular}{rrr}
    \toprule
    Sleep plateau & $V_{DD}$ energy & Control energy \\
    (ns) & (fJ/PE) & (fJ/PE) \\
    \midrule
       5 &   3.4506 & 13.9238 \\
      50 &  15.5239 & 13.9441 \\
     150 &  42.2864 & 13.9859 \\
     175 &  48.2865 & 13.9975 \\
     500 & 131.1477 & 14.2638 \\
    1,000 & 249.8182 & 14.5916 \\
   15,000 & 2,395.7000 & 19.6069 \\
    \bottomrule
  \end{tabular}
\end{table}

The electrical break-even duration is evaluated against leaving the PE ON
for the requested sleep plateau plus the measured 5.1-ns recovery window.  If
$T$ is the plateau duration, the per-PE gross benefit is
\begin{equation}
  E_{\mathrm{save}}(T)=P_{\mathrm{ON}}(T+5.1\,\mathrm{ns})
  -\left[E_{VDD}(T)+E_{\mathrm{ctrl}}(T)\right].
  \label{eq:hspice-break-even}
\end{equation}
The sign change lies between the directly simulated 150- and 175-ns points;
linear interpolation gives 165.3~ns.  At 1~GHz, the RTL conservatively rounds
this value upward to 166 cycles.  The additional one-cycle wake guard and
three certificate-pipeline cycles are controller-safety margins and are not
part of the 165.3-ns electrical break-even measurement.

\subsection{Connected-Block Validation and 128$\times$128 Scaling}
\label{sec:sa-scaling}

We test small-block scaling using directly connected 1$\times$1 and
4$\times$4 mapped arrays under the same ASAP7 TT device models.  The mapped
loads contain 2,225 and 33,858 standard-cell instances, respectively.  At the
documented $M=83$-per-PE sizing point, the 4$\times$4 header multiplicity is
1,326.  Table~\ref{tab:connected-scaling} compares the directly measured
ratios with the 16$\times$ PE-count ratio.

\begin{table}[t]
  \centering
  \caption{Direct 1$\times$1-to-4$\times$4 connected-block scaling.}
  \label{tab:connected-scaling}
  \begin{tabular}{lr}
    \toprule
    Metric & Measured 4$\times$4/1$\times$1 ratio \\
    \midrule
    Steady OFF power & 15.8568 \\
    Peak wake current & 15.7808 \\
    Entry total energy & 16.6646 \\
    Wake total energy & 14.9319 \\
    Entry latency & 1.04414 \\
    Wake latency & 0.93590 \\
    \bottomrule
  \end{tabular}
\end{table}

The near-16$\times$ power, current, and transition-energy ratios support a
linear per-PE energy model for the studied connected blocks; the latency
ratios show that latency is not multiplied by PE count.  The retained final
energy artifact also records direct 1$\times$1/4$\times$4 validation of the
$M=300$-per-PE droop assumption, but it does not preserve numerical 4$\times$4
$M=300$ measurements.  We therefore use the $M=300$ duration curve for energy
calculation while avoiding a quantitative full-array droop claim.

For workload evaluation, per-PE ON power, OFF power, and duration-dependent
event energy are multiplied by 16,384, the PE count of one 128$\times$128
array.  This is a linear scaling assumption:
\begin{equation}
  E_{128\times128}(T)=16{,}384\,E_{1\times1}(T).
  \label{eq:per-pe-energy-scaling}
\end{equation}
Because both the ON reference and event energy use the same factor, the
per-PE break-even duration is retained.  This extrapolation excludes
distributed power-grid resistance, header placement, isolation cells,
clock-control logic, architectural reset or weight-reload energy, and
full-array inrush interaction.

Controller area uses a separate structural extrapolation from synthesized
16$\times$16 and 32$\times$32 instances of the same two-stage PE RTL.  The
measured cell counts are 533,954 and 2,127,660, while the measured areas are
59,614.1937 and 237,764.4435~$\mu\mathrm{m}^{2}$.  Relative to four copies of
the 16$\times$16 design, the 32$\times$32 cell count and area differ by
$-0.3819$\% and $-0.2903$\%, respectively.  With PE count as the independent
variable, the frozen two-point extrapolation is
\begin{equation}
  Y_{128}=Y_{32}+\frac{16{,}384-1{,}024}{1{,}024-256}
  (Y_{32}-Y_{16})=21Y_{32}-20Y_{16}.
  \label{eq:structural-extrapolation}
\end{equation}
It yields 34,001,780 cells and 3.800769~mm$^{2}$ for a 128$\times$128 array.
These are extrapolated synthesis-stage values, not a measured full-array
place-and-route result.  A full 128$\times$128 implementation would be needed
to quantify routed timing, congestion, clock distribution, PDN droop, and
physical power-switch overhead.
